\section{Consequences}\label{sec:consequences}

Theorem~\ref{thm:main} allows us to apply two companion results for
nonamenable groups. The first concerns representations that are uniformly
close to being isometric but cannot be made unitary by an equivalent
Hilbert norm. The second concerns percolation on the simple Cayley graphs of \(F\)
associated with finite symmetric generating sets.
Neither companion result is used in the proof of nonamenability.

\subsection{Uniformly bounded representations}

\begin{corollary}\label{cor:nonunitarizable}
For every \(\varepsilon>0\), there are a separable complex Hilbert space
\(H\) and a representation \(\pi:F\to\mathrm{GL}(H)\) such that
\(\sup_{g\in F}\norm{\pi(g)}\le1+\varepsilon\), but no bounded invertible
operator \(S\) on \(H\) makes every \(S\pi(g)S^{-1}\) unitary.
\end{corollary}

\begin{proof}
The group \(F\) is countable, by its finite dyadic piecewise linear
description, and nonamenable by Theorem~\ref{thm:main}, so the conclusion
follows from the companion unitarizability theorem
\cite[Theorem~1.1]{OpenAIUnitarizability2026}.
\end{proof}

\subsection{Percolation on Cayley graphs}

For a finite symmetric generating set \(S\subset F\setminus\{\id\}\), let
\(G_S\) be the simple undirected Cayley graph with edges \(\{x,xs\}\).
In Bernoulli bond percolation on this fixed graph, each edge is independently
open with probability \(p\). Write \(p_c\) for the threshold for existence
of an infinite open cluster and \(p_u\) for the threshold for almost-sure
uniqueness of the infinite open cluster. The connection kernel
\(T_p(x,y)=\mathbb P_p(x\leftrightarrow y)\) acts on counting-measure
\(\ell^2(F)\), initially on finitely supported functions by summation over
\(y\); put \(p_{2\to2}=\sup\{p\in[0,1]:T_p\text{ is bounded on }\ell^2(F)\}\).

\begin{corollary}\label{cor:percolation}
For every such generating set \(S\),
\[
  \norm{T_{p_c}}_{2\to2}<\infty,\qquad p_c<p_{2\to2}\le p_u.
\]
For each fixed \(p\in(p_c,p_u)\), there are almost surely infinitely many
infinite open clusters. Moreover, there are deterministic
\(p_c<p_1<p_2<1\), depending on \(S\), such that in the coupling with
independent uniform \([0,1]\) edge labels \(U_e\), opening \(e\) when \(U_e\le p\),
there are almost surely infinitely many infinite open clusters
simultaneously for every \(p\in[p_1,p_2]\).
\end{corollary}

\begin{proof}
Apply the companion percolation result
\cite[Corollary~1.2]{OpenAIPercolation2026} to \(F\), which is nonamenable
by Theorem~\ref{thm:main}.
\end{proof}
